\uuid{mTCd}
\titre{Équation différentielle causale et transformée de Fourier}

\niveau{M1}
\module{Analyse}
\chapitre{Distributions et transformée de Fourier}
\sousChapitre{Équations différentielles dans $\mathcal{D}'_+$}
\theme{Solution fondamentale causale de l'opérateur $\partial_t+1$}
\auteur{Q. Liard}
\datecreate{ 2026-09-04}
\organisation{CentraleSupélec}
\difficulte{3}

\contenu{

\texte{
On cherche une distribution $G$ à support dans $\mathbb{R}_+$ vérifiant
$$
G'+G=\delta.
$$
On note $H$ la fonction de Heaviside. Dans tout l'exercice, on adopte la convention
$$
\widehat{u}(\xi)=\int_{\mathbb{R}}u(t)e^{-i\xi t}\,dt.
$$
Lorsque la formule d'inversion est applicable, elle s'écrit
$$
u(t)=\frac{1}{2\pi}\int_{\mathbb{R}}\widehat{u}(\xi)e^{i\xi t}\,d\xi.
$$
}

\begin{enumerate}

\item \question{Résoudre le problème de Cauchy
$$
\begin{cases}
z'+z=0,\\
z(0)=1.
\end{cases}
$$}

\reponse{
L'équation homogène $z'+z=0$ a pour solutions
$$
z(t)=Ce^{-t},
$$
où $C\in\mathbb{R}$. La condition initiale $z(0)=1$ impose $C=1$. Ainsi,
$$
\boxed{z(t)=e^{-t}}.
$$
}

\item \question{Soit $z$ la solution précédente. Vérifier que la distribution $Hz$ est une solution de
$$
G'+G=\delta.
$$}
\reponse{
Pour toute fonction $z$ de classe $C^1$, on a au sens des distributions
$$
(Hz)'=Hz'+z(0)\delta.
$$
Ici, $z(0)=1$ et $z'+z=0$. Par conséquent,
$$
(Hz)'+Hz
=
H(z'+z)+z(0)\delta
=
\delta.
$$
La distribution
$$
\boxed{G(t)=H(t)e^{-t}}
$$
est donc une solution de $G'+G=\delta$. Son support est contenu dans $\mathbb{R}_+$.
}

\item \question{Avec la convention de Fourier donnée dans l'énoncé, rappeler les expressions de $\widehat{G'}$ et de $\widehat{\delta}$, puis transformer l'équation
$$
G'+G=\delta.
$$
En déduire l'expression de $\widehat{G}$.}

\reponse{
Avec la convention
$$
\widehat{u}(\xi)=\int_{\mathbb{R}}u(t)e^{-i\xi t}\,dt,
$$
la dérivation dans la variable $t$ devient une multiplication par $i\xi$ :
$$
\widehat{G'}(\xi)=i\xi\,\widehat{G}(\xi).
$$
De plus,
$$
\widehat{\delta}(\xi)=1.
$$
La transformation de Fourier de l'équation donne alors
$$
i\xi\,\widehat{G}(\xi)+\widehat{G}(\xi)=1,
$$
soit
$$
(1+i\xi)\widehat{G}(\xi)=1.
$$
Comme $1+i\xi\neq 0$ pour tout $\xi\in\mathbb{R}$, on obtient
$$
\boxed{\widehat{G}(\xi)=\frac{1}{1+i\xi}}.
$$
}

\item \question{Vérifier directement que la transformée de Fourier de la fonction
$$
t\longmapsto H(t)e^{-t}
$$
est égale à $\dfrac{1}{1+i\xi}$.}

\reponse{
La fonction $t\mapsto H(t)e^{-t}$ appartient à $L^1(\mathbb{R})$. Pour tout $\xi\in\mathbb{R}$,
$$
\widehat{H e^{-\cdot}}(\xi)
=
\int_{\mathbb{R}}H(t)e^{-t}e^{-i\xi t}\,dt
=
\int_0^{+\infty}e^{-(1+i\xi)t}\,dt.
$$
Comme la partie réelle de $1+i\xi$ vaut $1>0$, l'intégrale converge et
$$
\widehat{H e^{-\cdot}}(\xi)
=
\left[
-\frac{e^{-(1+i\xi)t}}{1+i\xi}
\right]_{0}^{+\infty}.
$$
Le terme exponentiel tend vers $0$ lorsque $t\to+\infty$. Ainsi,
$$
\boxed{
\widehat{H e^{-\cdot}}(\xi)=\frac{1}{1+i\xi}
}.
$$
Ce résultat coïncide avec celui obtenu en transformant directement l'équation différentielle.
}

\item \question{Plus généralement, soit $a>0$. Déterminer une solution $G_a\in\mathcal{D}'_+$ de
$$
G_a'+aG_a=\delta,
$$
puis calculer sa transformée de Fourier avec la même convention.}

\reponse{
On résout d'abord
$$
z_a'+az_a=0,
\qquad
z_a(0)=1.
$$
On obtient
$$
z_a(t)=e^{-at}.
$$
Par la même formule de dérivation,
$$
(Hz_a)'+aHz_a
=
H(z_a'+az_a)+z_a(0)\delta
=
\delta.
$$
On peut donc choisir
$$
\boxed{G_a(t)=H(t)e^{-at}}.
$$
En transformant l'équation,
$$
i\xi\,\widehat{G_a}(\xi)+a\widehat{G_a}(\xi)=1.
$$
Par conséquent,
$$
\boxed{\widehat{G_a}(\xi)=\frac{1}{a+i\xi}}.
$$
On retrouve également cette formule par le calcul direct
$$
\int_0^{+\infty}e^{-(a+i\xi)t}\,dt=\frac{1}{a+i\xi}.
$$
}

\end{enumerate}

}
